% !TeX root = ../main.tex
\section{问题分析}
\subsection{问题一的分析}
光伏出力主要集中在白天，而负荷高峰与高电价时段并不完全重合，说明光伏发电与用电需求之间存在一定的时序错配。
同时，充放电过程还存在能量损耗。

由于题目给出的各项数据均为确定值，且每天采用相同的变化曲线，可将一天划分为连续的若干时段，并将问题视为一个单日确定性调度问题。
据此，可以以全天购电费用最小为主要目标，建立单日储能调度线性规划模型，并进一步采用两阶段优化思路对调度方案进行调整。

\subsection{问题二的分析}
本问的难点在于 0:00 制定计划时实际负荷与光伏出力尚未确定，购电偏差可能造成计划电量闲置或高价紧急购电的情况出现。
因此，可以采用季节朴素、岭回归和梯度提升树预测次日负荷与光伏，并以加权绝对百分比误差评价预测效果。
在此基础上，建立引入条件风险价值（CVaR）的两阶段随机线性规划模型，兼顾购电成本与预测偏差风险。

\subsection{问题三的分析}
【说明求解路径、模型之间的联系与结果检验方式。】
\subsection{问题四的分析}
【说明问题四的目标、难点、建模方法及与前三问的联系。】
\subsection{整体研究思路}
【在下图位置放入总体流程图，并用文字解释各阶段的输入和输出。】
\begin{figure}[htbp]
  \centering
  % 将下行替换为：\includegraphics[width=0.85\linewidth]{figures/framework.pdf}
  \figureplaceholder{总体研究流程图}{数据整理 $\rightarrow$ 模型建立 $\rightarrow$ 求解检验 $\rightarrow$ 结论}
  \caption{整体研究思路}\label{fig:framework}
\end{figure}
